\section*{Замечания к исходным рабочим заметкам}
Решения упражнений \S\,3 основаны на заметках из каталога
\texttt{notes/claude/} (файлы \texttt{p21-*.md}); в этом документе они
переписаны так, чтобы каждое решение было самостоятельным: решение
упражнения~1 содержит полные доказательства леммы (L) и свойств таблиц
истинности, остальные решения напоминают нужные факты со ссылкой на него.
По сравнению с заметками исправлено следующее.
\begin{itemize}
\item В разминочной таблице заметки к упр.~2 формула
  $(A\cap B)\cup(\ov A\cap I)$ отмечена как самодвойственная (<<обе
  $\equiv\ov A\cup B$>>). Это неверно: $F\equiv\ov A\cup B$, а
  $F^*=(A\cup B)\cap\ov A\equiv\ov A\cap B$ (что правильно вычислено ниже
  в той же заметке при проверке леммы (L)). В книге эта пара приведена лишь
  как пример \emph{двойственных} формул.
\item В заметке к упр.~3 сказано, что для произвольной формулы <<ни из
  $F\cap F^*=\emptyset$ не следует $F\cup F^*=I$, ни наоборот>>. На самом
  деле эти условия равносильны для любой формулы (оба равносильны
  четности); замечание удалено.
\item Универсальное множество в примере $\{0,1\}^n$ обозначено $U$, чтобы не
  путать его с константой $I$.
\end{itemize}
Все численные утверждения решений проверены программами из каталога
\texttt{exercises/artifacts/checks/} (результаты "--- в файле
\texttt{RESULTS.txt}).
